% Generated by render_live_results.py; edit the source evidence only via a new run.
本次從固定 \code{3cbaa518f0dc} 的乾淨 detached worktree 重建 runtime，沿用原 engine，不重建模型。另行編譯僅增加 host enqueue／receive timestamps 的 benchmark；runtime 原檔未修改。測試為 HR、G1、Graph ON、NSA／量測框 OFF、result ring=4，使用 MOT17-10-SDP 前 140 幀，每 session 暖機 20 幀、量測 120 幀。Single 與 queued 交替各四次，共八個 fresh sessions、每模式 480 筆暖機後樣本。\cite{live}

\begin{table}[H]
\caption{本次 G1 提交至結果交付延遲，單位 ms}
\small
\noindent\begin{tabularx}{\textwidth}{@{}Yrr@{}}
\toprule
統計 & Single：同幀立即 receive & Queued：四格 ring \\
\midrule
平均 & 6.712865 & 20.254368 \\
p50 & 6.680952 & 20.478764 \\
p95 & 6.981673 & 21.181116 \\
p99 & 7.228440 & 21.571266 \\
最小值 & 6.392425 & 12.934183 \\
最大值 & 9.288572 & 24.747265 \\
\bottomrule
\end{tabularx}
\end{table}

Queued wall interval 平均為 \textbf{5.138200 ms/frame}，換算吞吐為 \textbf{194.621 FPS}。它與 queued 交付平均 20.254 ms、p95 21.181 ms 是不同指標；較深佇列的吞吐不能代替同幀交付時間。

\begin{table}[H]
\caption{本次初始化與首幀時間：八個 fresh sessions，單位 ms}
\small
\noindent\begin{tabularx}{\textwidth}{@{}Yrrr@{}}
\toprule
項目 & 最小值 & 中位數 & 最大值 \\
\midrule
Session constructor & 435.330 & 443.964 & 3371.001 \\
首幀 enqueue 至 receive & 16.445 & 16.895 & 23.062 \\
\bottomrule
\end{tabularx}
\end{table}

\textbf{一致性與環境。} 八份 140 幀完整 packet stream 逐 byte 相同，包含暖機、未用 rows 與 reserved bytes；source、engine、輸入及 binaries 的量測前後 hashes 均相同。Driver 為 621.14，CUDA 13.4.92、TensorRT 10.16.1.11；與歷史 driver 617.42 不同，不合併樣本、不宣稱跨時段加速。封包相同指本次 sessions 之間的一致性，不是新一次品質驗收。

\textbf{遙測與適用限制。} 每秒一次的 12 筆程序期遙測包含 constructor 與暖機，GPU utilization 平均 46.5\%、最大 100\%，最高溫 68\,°C、功率 139.46 W、顯存 495 MiB；不能當作穩態利用率。只有一序列 prefix 與每模式四次 sessions，JPEG 已預載、未鎖 clocks；每對固定 single 先、queued 後，順序未反平衡。Host timestamp 會有小量成本，所有長尾保留。

\textbf{統計口徑。} 表中 p50／p95／p99 採 nearest-rank，初始化表的中位數則為中間兩值平均。權威逐幀樣本為 \code{*.delivery.csv}，終點是 receive 返回、packet 檢查與複製之前。原 stock harness 的 \code{runs.csv} 與 single CSV 另含後續檢查／複製，並保留原樣；本表不混用兩種口徑。
